\begin{textsample}{POS Dim 2 – sp – Score 17.00 – sp/sp\_em\_44.txt}  \label{ex:f2_pos_235}
\textbf{Itinerário} Formativo - Formação \textbf{Técnica} e Profissional A etapa do Ensino Médio objetiva a construção de novos saberes e habilidades para a solução de problemas do mundo real, mobilizando linguagens, raciocínio lógico-matemático, conhecimentos sócio históricos, científicos, tecnológicos, artísticos e culturais no contexto da sociedade contemporânea, além de competências pessoais como capacidade de trabalhar de modo autônomo e colaborativo, respeitando valores e atitudes éticas e ambientais.
Além disso, o estudante do Ensino Médio deverá ser capaz de argumentar com base em fatos e fontes diversas, cuidar de sua saúde física e emocional e agir com flexibilidade nos campos de atuação social.
A formação \textbf{técnica} e profissional, assim como as áreas do conhecimento, compõe a flexibilização da etapa do Ensino Médio, nos termos da Lei nº 13.415/2017, proporcionando ao estudante desenvolver e fortalecer a sua autonomia, considerando seu projeto de vida, permitindo ter uma visão ampla do mundo e tomar decisões em sua vida escolar, profissional e pessoal.
O parágrafo 2º do \textbf{Artigo} 12 das DCNEM, nos termos da Resolução nº 3, de 21 de \textbf{novembro} de 2018, \textbf{reiterados} na \textbf{Portaria} nº 1.432, de 28 de dezembro de 2018, descreve que os \textbf{itinerários} formativos das diferentes áreas e da formação \textbf{técnica} e profissional devem ser organizados considerando quatro eixos estruturantes.
Especificamente sobre o \textbf{itinerário} de formação \textbf{técnica} e profissional, o parágrafo 14 do mesmo \textbf{artigo} traz : “ [... ] deve observar a integralidade de ocupações \textbf{técnicas} reconhecidas pelo setor produtivo, tendo como referência a Classificação Brasileira de Ocupações ( CBO ) ”. ( BRASIL, 2018 ) Ressalta -se que as \textbf{Diretrizes} Curriculares Nacionais para a Educação Profissional \textbf{Técnica} de Nível Médio são definidas pela Resolução CNE/CEB nº 6, de 20 de setembro de 2012, regramento que reúne os “ [... ] princípios e critérios a serem observados pelos sistemas de ensino e pelas \textbf{instituições} de ensino públicas e privadas, na organização e no planejamento, desenvolvimento e avaliação ” da modalidade. ( BRASIL, 2012 ) Os pressupostos de implementação dos \textbf{itinerários} formativos apontados pela \textbf{Portaria} nº 1.432/2018 trazem orientações para que os quatro eixos estruturantes sejam incorporados visando garantir que o estudante experimente diversas situações de aprendizagem e tenha capacidade de desenvolver um amplo conjunto de habilidades.
Especificamente para a formação \textbf{técnica} e profissional, a \textbf{Portaria} traz : > No caso da Formação \textbf{Técnica} e Profissional, os \textbf{Itinerários} também se organizam a partir da integração dos diferentes eixos estruturantes, ainda que as habilidades a eles associadas somem -se a outras habilidades básicas requeridas indistintamente pelo mundo do trabalho e a habilidades específicas requeridas pelas distintas ocupações, conforme \textbf{previsto} no \textbf{Catálogo} Nacional de \textbf{Cursos} \textbf{Técnicos} - CNCT e na Classificação Brasileira de Ocupações - CBO. ( BRASIL, 2018 ) O \textbf{Currículo} Paulista etapa do Ensino Médio traz em seu escopo o \textbf{itinerário} \textbf{técnico} profissional, com um conjunto selecionado e coerente de saberes e capacidades, bem como a mobilização de competências socioemocionais e \textbf{técnicas} direcionadas à ação e à reflexão em uma área ou campo do saber – que se constitui, em última instância, no próprio \textbf{Currículo}.
Tal formação visa ao preparo para o mundo do trabalho em cargos, funções ou de modo autônomo, contribuindo para a inserção do cidadão na sociedade.
Segue um breve \textbf{detalhamento} das possibilidades de desenvolvimento da formação \textbf{técnica} e profissional em cada eixo estruturante : INVESTIGAÇÃO CIENTÍFICA A investigação de aspectos científicos e práticos tem o propósito de estimular no estudante o interesse em desenvolver pesquisas relacionadas às situações-problema do mundo do trabalho, contemplando amplas possibilidades de transformação de produtos, processos, serviços e outros elementos intrínsecos aos seus anseios profissionais.
Sugere -se que o estudante pesquise ambientes de acesso cotidiano, como escola e moradia, para vislumbrar aspectos do trabalho nesses ambientes, fazendo uma analogia com aquilo praticado nas organizações – relacionamento interpessoal, departamentalização, infraestrutura, entre várias alternativas.
PROCESSOS CRIATIVOS As possibilidades deste eixo estruturante contemplam o desenvolvimento da capacidade do estudante em aprender com base em processos inovadores, associados à economia criativa, geração de riqueza e sustentabilidade, de modo que seja estimulada a investigação científica de forma conjunta com essas ações.
Sugere -se que o estudante, a partir de leituras específicas e atuais, coloque em prática seu pensar e fazer criativos, e consequentemente desenvolva competências e habilidades voltadas à idealização e prototipação de soluções requeridas pela área profissional.
MEDIAÇÃO E INTERVENÇÃO SOCIOCULTURAL As ações deste eixo estruturante proporcionam condições para múltiplas aprendizagens, que podem ser embasadas em experiências socioculturais em ambientes voltados à área profissional/eixo tecnológico da formação \textbf{técnica}.
Os conhecimentos a serem mobilizados devem contemplar ações individuais e coletivas de mediação e intervenção sobre nossa sociedade e cultura, de modo a promover o engajamento e o protagonismo dos estudantes em situações que demandam atitudes e propostas de reflexão e de mudança – ações de conscientização, campanhas relacionadas ao respeito à diversidade, entre outras possibilidades.
EMPREENDEDORISMO A prática do empreendedorismo para a formação profissional, embora \textbf{pareça} elementar, deve ser fortalecida em amplas frentes, não se voltando apenas à implementação de negócios, mas também a ações intraempreendedoras e propostas que contemplem os anseios do estudante – seu projeto de vida é um exemplo prático disso.
Para desenvolver a temática deste eixo estruturante, pode -se adotar como premissa o desenvolvimento de competências transversais que mobilizem conhecimentos e busquem novas oportunidades de negócio e de atuação profissional, permitindo ao estudante resolver problemas relacionados ao mundo do trabalho e ao seu cotidiano.
Como exemplo, pode -se criar um portfólio ao longo do Ensino Médio, composto de trabalhos, pesquisas, produções práticas, protótipos, entre outras possibilidades.
A educação profissional é marcada pela combinação das demandas dos setores produtivos, dos interesses dos indivíduos e dos interesses coletivos, devendo, assim, preparar o cidadão para o desempenho de “ profissões ”, cada vez mais fluidas, intangíveis e mutantes.
O \textbf{trabalhador} deve estar habituado e preparado para a adaptação contínua às possíveis relações profissionais, às novas tecnologias e inovações, ao posicionamento intelectual, político e filosófico dos atores sociais, incluindo concepções e visões de mundo, comportamento, condutas e valores.
O \textbf{Currículo} Paulista, em sua parte \textbf{flexível}, apresenta em um dos \textbf{itinerários} formativos a formação \textbf{técnica} e profissional, que direciona o planejamento, a sistematização e o desenvolvimento de perfis profissionais, de atribuições, de atividades, de competências, de habilidades e de bases tecnológicas, valores e conhecimentos, organizados em componentes curriculares e por eixo tecnológico ou área de conhecimento.
O foco do ensino profissionalizante deve estar alinhado aos interesses do estudante.
Os conhecimentos ( temas relativos à vida contemporânea e ao cânone cultural de cada sociedade ), as habilidades, incluindo as inclinações \textbf{técnicas}, tecnológicas e científicas, deverão estar relacionados aos seus anseios e seu projeto de vida.
O \textbf{itinerário} formativo referente à formação \textbf{técnica} e profissional deve ser desenvolvido em um contexto contemporâneo, altamente volátil e que contemple a informatização e digitalização de processos de ensino e aprendizagem, além de materiais lúdicos, plataformas de ensino diferenciadas e gamificação de processos pedagógicos, mostrando -se sempre aderente às demandas locais.
A seguir encontra -se o organizador curricular do \textbf{itinerário} formativo da formação \textbf{técnica} e profissional, que contempla as habilidades relacionadas às competências gerais da Educação Básica, \textbf{destinadas} a desenvolver perfis profissionais requeridos pelo mundo do trabalho, e os pressupostos metodológicos que visam orientar e indicar possibilidades para a concretização das aprendizagens esperadas, para cada eixo estruturante.

% matched lemmas: artigo, catálogo, currículo, curso, destinar, detalhamento, diretriz, flexível, instituição, itinerário, novembro, parecer, portaria, prever, reiterar, trabalhador, técnico
\end{textsample}
